\section{Experiments}
\label{sec:eval}

\subsection{Experiment Setup}
\label{sec:setup}

We evaluate \sys on Qwen3-8B at temperature $0$ with prefix caching and tensor parallelism $1$, on four A100 40GB GPUs.
By default we use the ordered pair in Equation~\ref{eq:pair} and probe length $t_0{=}128$.
We evaluate with GSM8K under three decode-time hosts and with MATH-500 under the probe rule.
We also report peak KV, quality, and offered load on a serving stack whose live occupancy is the committed spine, using six model families.
Please refer to Appendix~\ref{app:details} for host parameters and metric definitions.
Note that \sys is complementary to the hosts, since our approach does not modify their scoring rules or require additional test-time model calls.

\paragraph{GSM8K hosts.}
Sixteen GSM8K problems.
Each problem is a fan-out of $k \in \{4,8,16\}$ residuals and a decode budget $D{=}512$.
Each fan-out is a controlled mix: half the residuals are live strategy instructions, and half alternate a repeated loop with illegal text.
Index $0$ is always a live strategy.
\emph{Any-finished} counts a problem if any admitted branch produces the reference number.
\emph{Winner} counts a problem if the branch selected as the winner does.
Both are reported over $16$.
ESC, speculative rejection, and DPTS are the decode hosts.
The host alone submits every residual.
\sys applies $\tau_{\mathrm{pre}}$ first and leaves the host's cut in place.

\paragraph{MATH-500 probe.}
The first $32$ MATH-500 problems at level $4$ or higher, $k \in \{4,8,16\}$, $D{=}1024$.
Half the residuals are contest strategies and half are loops or illegal text.
The decode bar is $\tau_{\mathrm{dec}}$.
Full fan-out decodes every residual to $D$.
Discard drops every residual below that bar before prefill.
Probe admits those residuals for $t_0$ tokens and extends only when the generated prefix clears $\tau_{\mathrm{dec}}$.

\subsection{Result Analysis}
\label{sec:analysis}

We summarize the results in Table~\ref{tab:hosts} and Figure~\ref{fig:results}.
\sys significantly reduces the prefill shared by three decode-time hosts, while maintaining any-finished accuracy.

\begin{table}[t]
\centering
\small
\setlength{\tabcolsep}{4pt}
\begin{tabular}{@{}lrrrr@{}}
\toprule
$k{=}16$ & Prefill & Host & $+\tau_{\mathrm{pre}}$ & Any \\
\midrule
ESC & $5008 \to 3848$ & 17.05\,s & \textbf{14.13\,s} & $16/16$ \\
SR & $5008 \to 3848$ & 13.14\,s & \textbf{12.47\,s} & $16/16$ \\
DPTS & $5008 \to 3848$ & 11.86\,s & \textbf{11.49\,s} & $16/16$ \\
\bottomrule
\end{tabular}
\caption{Prefill pruning at $\tau_{\mathrm{pre}}$ in front of three hosts.
Any-finished stays $16/16$.
The prefill column is identical across hosts because the gate, not the host cut, changes $W$, as predicted by Proposition~\ref{prop:invariance}.}
\label{tab:hosts}
\end{table}

Specifically, with the prefill gate off, the three hosts share one prefill of $5008$ tokens.
Turning on $\tau_{\mathrm{pre}}$ drops the loops, leaves illegal text, and cuts that prefill to $3848$ on every host.
ESC moves the most, because its baseline decodes the loops out to $D$.
Speculative rejection and DPTS move less, because they were already stopping the loop and illegal half at their own cut.
The gate removes the prefill of the loop, which none of those cuts avoids.
This highlights a deficiency of pruning only after tokens exist: the host can stop an abortable residual, but it cannot unspend $W$.

\begin{figure*}[t]
\centering
\resizebox{\textwidth}{!}{\begin{tikzpicture}[
  x=1cm, y=1cm,
  font=\sffamily\scriptsize,
  >=Latex,
  line cap=round,
]
% ---------- (a) grouped bars, 0.18 cm per second ----------
\node[font=\scriptsize\bfseries\sffamily, anchor=west] at (0.05,4.72) {(a) End-to-end time, $k{=}16$};
\begin{scope}[shift={(0.70,0.62)}]
  \draw[fsgray!35] (0,0) -- (6.70,0);
  \draw[fsgray!35] (0,0) -- (0,3.42);
  \foreach \s/\lab in {0/0, 5/5, 10/10, 15/15} {
    \draw[fsgray!22] (0,{0.18*\s}) -- (6.70,{0.18*\s});
    \node[font=\tiny\sffamily, anchor=east, text=fsgray] at (-0.08,{0.18*\s}) {\lab};
  }
  \node[font=\tiny\sffamily, rotate=90, anchor=south, text=fsgray] at (-0.62,1.55) {seconds};

  \fill[fsgray!50] (0.28,0) rectangle (0.78,{0.18*17.05});
  \fill[fsblue!60] (0.82,0) rectangle (1.32,{0.18*14.13});
  \fill[fsgreen!60] (1.36,0) rectangle (1.86,{0.18*10.92});
  \node[font=\tiny\sffamily, anchor=south] at (0.53,{0.18*17.05+0.04}) {17.1};
  \node[font=\tiny\sffamily, anchor=south] at (1.07,{0.18*14.13+0.04}) {14.1};
  \node[font=\tiny\sffamily, anchor=south] at (1.61,{0.18*10.92+0.04}) {10.9};

  \fill[fsgray!50] (2.48,0) rectangle (2.98,{0.18*13.14});
  \fill[fsblue!60] (3.02,0) rectangle (3.52,{0.18*12.47});
  \fill[fsgreen!60] (3.56,0) rectangle (4.06,{0.18*10.80});
  \node[font=\tiny\sffamily, anchor=south] at (2.73,{0.18*13.14+0.04}) {13.1};
  \node[font=\tiny\sffamily, anchor=south] at (3.27,{0.18*12.47+0.04}) {12.5};
  \node[font=\tiny\sffamily, anchor=south] at (3.81,{0.18*10.80+0.04}) {10.8};

  \fill[fsgray!50] (4.68,0) rectangle (5.18,{0.18*11.86});
  \fill[fsblue!60] (5.22,0) rectangle (5.72,{0.18*11.49});
  \fill[fsgreen!60] (5.76,0) rectangle (6.26,{0.18*10.90});
  \node[font=\tiny\sffamily, anchor=south] at (4.93,{0.18*11.86+0.04}) {11.9};
  \node[font=\tiny\sffamily, anchor=south] at (5.47,{0.18*11.49+0.04}) {11.5};
  \node[font=\tiny\sffamily, anchor=south] at (6.01,{0.18*10.90+0.04}) {10.9};

  \node[font=\tiny\sffamily] at (1.07,-0.28) {ESC};
  \node[font=\tiny\sffamily] at (3.27,-0.28) {SR};
  \node[font=\tiny\sffamily] at (5.47,-0.28) {DPTS};
\end{scope}

\fill[fsgray!50] (0.18,4.28) rectangle (0.48,4.48);
\node[font=\tiny\sffamily, anchor=west] at (0.54,4.38) {host};
\fill[fsblue!60] (1.42,4.28) rectangle (1.72,4.48);
\node[font=\tiny\sffamily, anchor=west] at (1.78,4.38) {$+\tau_{\mathrm{pre}}$};
\fill[fsgreen!60] (2.78,4.28) rectangle (3.08,4.48);
\node[font=\tiny\sffamily, anchor=west] at (3.14,4.38) {prefill at $\tau_{\mathrm{dec}}$};

% ---------- (b) MATH-500 ----------
\begin{scope}[shift={(8.35,0.62)}]
  \node[font=\scriptsize\bfseries\sffamily, anchor=west] at (-0.55,4.10) {(b) MATH-500, any-finished};
  \fill[fsgray] (0.05,3.12) rectangle (0.28,3.28);
  \node[font=\tiny\sffamily, anchor=west] at (0.34,3.20) {full};
  \fill[fsred] (0.05,2.84) rectangle (0.28,3.00);
  \node[font=\tiny\sffamily, anchor=west] at (0.34,2.92) {discard};
  \fill[fsgain] (0.05,2.56) rectangle (0.28,2.72);
  \node[font=\tiny\sffamily, anchor=west] at (0.34,2.64) {probe};

  \draw[fsgray!35] (0,0) -- (5.20,0);
  \draw[fsgray!35] (0,0) -- (0,3.35);
  \foreach \s/\lab in {0/0, 5/5, 10/10, 15/15, 20/20} {
    \draw[fsgray!18] (0,{0.155*\s}) -- (5.20,{0.155*\s});
    \node[font=\tiny\sffamily, anchor=east, text=fsgray] at (-0.08,{0.155*\s}) {\lab};
  }
  \foreach \s/\lab in {0/0, 25/25, 50/50, 75/75, 100/100} {
    \draw[fsgray!40] ({0.046*\s},0) -- ({0.046*\s},-0.07);
    \node[font=\tiny\sffamily, anchor=north, text=fsgray] at ({0.046*\s},-0.10) {\lab};
  }
  \node[font=\tiny\sffamily, text=fsgray, anchor=north] at (2.55,-0.46) {seconds};
  \node[font=\tiny\sffamily, rotate=90, anchor=south, text=fsgray] at (-0.68,1.40) {correct / 32};

  \draw[fsgray, line width=0.85pt] ({0.046*25.3},{0.155*12}) -- ({0.046*52.5},{0.155*15}) -- ({0.046*104.1},{0.155*20});
  \foreach \t/\a in {25.3/12, 52.5/15, 104.1/20} {
    \fill[fsgray] ({0.046*\t},{0.155*\a}) circle (1.55pt);
  }
  \draw[fsred, line width=0.85pt] ({0.046*19.1},{0.155*6}) -- ({0.046*25.6},{0.155*11}) -- ({0.046*52.1},{0.155*13});
  \foreach \t/\a in {19.1/6, 25.6/11, 52.1/13} {
    \fill[fsred] ({0.046*\t},{0.155*\a}) circle (1.55pt);
  }
  \draw[fsgain, line width=0.95pt] ({0.046*24.7},{0.155*9}) -- ({0.046*43.3},{0.155*13}) -- ({0.046*72.2},{0.155*20});
  \foreach \t/\a in {24.7/9, 43.3/13, 72.2/20} {
    \fill[fsgain] ({0.046*\t},{0.155*\a}) circle (1.7pt);
  }
  \node[font=\tiny\sffamily, anchor=south west] at ({0.046*72.2+0.08},{0.155*20+0.02}) {$k{=}16$};
\end{scope}
\end{tikzpicture}
}
\caption{Qwen3-8B.
(a)~GSM8K controlled mix, $k{=}16$, $n{=}16$, $D{=}512$.
Each group is one host.
The gray bar is the host alone, blue adds $\tau_{\mathrm{pre}}$, and green sets the prefill bar to $\tau_{\mathrm{dec}}$.
Any-finished accuracy holds on every bar.
At the higher prefill bar the three hosts meet.
(b)~MATH-500, $n{=}32$, $D{=}1024$.
Each polyline connects $k{=}4$, $k{=}8$, and $k{=}16$.
Discarding residuals below $\tau_{\mathrm{dec}}$ before prefill is fast and less accurate.
The probe under $(\tau_{\mathrm{pre}}, \tau_{\mathrm{dec}})$ meets the full fan-out at $k{=}16$.}
\label{fig:results}
\end{figure*}

Figure~\ref{fig:results}(a) raises the prefill bar from $\tau_{\mathrm{pre}}$ to $\tau_{\mathrm{dec}}$, which also drops illegal text.
Any-finished holds at $k{=}16$.
End-to-end time converges across the three hosts, and the host's later cut fires on nobody: the prefill gate has already removed every residual that cut would have rejected.

\begin{table*}[t]
\centering
\small
\setlength{\tabcolsep}{5pt}
\begin{tabular}{@{}lrrrrrrr@{}}
\toprule
Setting & Admit & Cut later & Prefill & Decode & Time & Any & Winner \\
\midrule
\multicolumn{8}{@{}l}{\emph{Decode pruning only.} No $\tau_{\mathrm{pre}}$. Every residual is prefilled.} \\
ESC, budget $D{=}512$ & 16 & 0 & 5008 & 131072 & 17.05\,s & $16/16$ & $11/16$ \\
SR, horizon $256$ & 16 & 8 & 5008 & 98304 & 13.14\,s & $16/16$ & $12/16$ \\
DPTS, mini-step $100$ & 16 & 8 & 5008 & 78336 & 11.86\,s & $16/16$ & $12/16$ \\
\midrule
\multicolumn{8}{@{}l}{\emph{Prefill pruning only.} Decode runs each admitted residual to $D{=}512$.} \\
$\tau_{\mathrm{pre}}$ & 12 & 0 & 3848 & 98304 & 14.13\,s & $16/16$ & $11/16$ \\
$\tau_{\mathrm{dec}}$ as prefill & 8 & 0 & 2768 & 65536 & \textbf{10.92\,s} & $16/16$ & $10/16$ \\
\midrule
\multicolumn{8}{@{}l}{\emph{Both.} The same prefill bar, then the host cut on what remains.} \\
SR $+\tau_{\mathrm{pre}}$ & 12 & 4 & 3848 & 81920 & 12.47\,s & $16/16$ & $10/16$ \\
DPTS $+\tau_{\mathrm{pre}}$ & 12 & 4 & 3848 & 71936 & 11.49\,s & $16/16$ & $12/16$ \\
SR $+\tau_{\mathrm{dec}}$ & 8 & 0 & 2768 & 65536 & \textbf{10.80\,s} & $16/16$ & $10/16$ \\
DPTS $+\tau_{\mathrm{dec}}$ & 8 & 0 & 2768 & 65536 & \textbf{10.90\,s} & $16/16$ & $10/16$ \\
\bottomrule
\end{tabular}
\caption{Ablation at $k{=}16$ on the GSM8K controlled mix.
``Cut later'' is the number of admitted residuals the host stops at its own horizon.
Any-finished is unchanged.
Winner accuracy moves by at most two problems.
When the prefill bar is $\tau_{\mathrm{dec}}$, the host cut fires on nobody and the decode columns meet.}
\label{tab:ablation}
\end{table*}

Table~\ref{tab:ablation} separates the two gates.
Decode pruning alone never changes prefill, and the time ranking is the ranking of the cut.
Prefill pruning alone, with no host cut, is already faster than ESC at $\tau_{\mathrm{pre}}$ and faster than all three hosts when the prefill bar is $\tau_{\mathrm{dec}}$.
Stacking a host on $\tau_{\mathrm{pre}}$ still stops some of the admitted residuals.
Stacking a host on $\tau_{\mathrm{dec}}$ does not, because the admitted set is the live half.
Any-finished holds throughout.
Winner accuracy is not identical to any-finished: dropping illegal residuals can change which branch the host ranks first.
The gate preserves a correct branch in the admitted set.
Appendix~\ref{app:width} reports $k{=}4$ and $k{=}8$, where the higher prefill bar does drop any-finished accuracy.

The two scores of a loop or an illegal residual are not interchangeable, which is the empirical content of the score gap in Section~\ref{sec:gap}.
On the residual, illegal text lies above $\tau_{\mathrm{pre}}$.
On a generated prefix that is still illegal, the decode score lies below $\tau_{\mathrm{dec}}$.
Applying $\tau_{\mathrm{dec}}$ as a prefill bar therefore drops illegal text, which \eqref{eq:keep} at $\tau_{\mathrm{pre}}$ would have admitted.
On MATH-500 at large $k$ that is the discard row of Table~\ref{tab:probe}.
Applying $\tau_{\mathrm{pre}}$ as a decode bar would do the opposite: an illegal prefix still clears $\tau_{\mathrm{pre}}$, so the probe would extend a residual the decode rule is designed to stop.

\begin{table}[t]
\centering
\small
\setlength{\tabcolsep}{3.5pt}
\begin{tabular}{@{}lccc@{}}
\toprule
 & $k{=}4$ & $k{=}8$ & $k{=}16$ \\
\midrule
Full & 12/32, 25.3\,s & 15/32, 52.5\,s & 20/32, 104.1\,s \\
Discard & 6/32, 19.1\,s & 11/32, 25.6\,s & 13/32, 52.1\,s \\
Probe & 9/32, 24.7\,s & 13/32, 43.3\,s & \textbf{20/32, 72.2\,s} \\
\bottomrule
\end{tabular}
\caption{MATH-500, Qwen3-8B, $n{=}32$, $D{=}1024$.
Each cell is any-finished accuracy and end-to-end time.
At $k{=}16$ the probe matches the full fan-out and is $31\%$ faster.
At $k{=}4$ and $k{=}8$ it does not fully recover accuracy.}
\label{tab:probe}
\end{table}

Table~\ref{tab:probe} and Figure~\ref{fig:results}(b) compare the three prefill policies on MATH-500.
Discarding every residual below $\tau_{\mathrm{dec}}$ is the fastest curve and the least accurate.
The probe spends $t_0$ tokens on the uncertain residuals, then extends a residual only when its prefix clears $\tau_{\mathrm{dec}}$.
At $k{=}16$ any-finished returns to the full fan-out.
At $k{=}4$ and $k{=}8$ it does not.
Recovery is a wide-fan-out result.
Appendix~\ref{app:probe} sweeps $t_0 \in \{64,128\}$, loop-skipping, and a decode-only short budget.
The $k{=}16$ match to full accuracy is specific to $t_0{=}128$ among the probe lengths we compare; $t_0{=}64$ reaches $18/32$.

The prefill gate does not lengthen the committed spine.
Appendix~\ref{app:spine} measures that spine against recomputation and automatic prefix caching.
Peak KV stays below the cache on six model families, because live occupancy is the spine.
On Qwen3-8B the pruned system keeps GSM8K quality with that cache and lowers end-to-end time.
Under a $1$\,s P99 time-to-first-token line, the shorter spine admits a higher offered load.

\begin{table}[t]
\centering
\small
\setlength{\tabcolsep}{4pt}
\begin{tabular}{@{}lrrrr@{}}
\toprule
$k{=}16$, bar $\tau_{\mathrm{dec}}$ & Admit & Prefill & ESC & SR \\
\midrule
Keep one, ignore score & 1 & 414 & $11/16$ & $9/16$ \\
Keep all above $\tau_{\mathrm{dec}}$ & 8 & 2768 & $\mathbf{16/16}$ & $\mathbf{16/16}$ \\
\bottomrule
\end{tabular}
\caption{Winner-only admission on the GSM8K mix.
A single reserved slot drops any-finished accuracy on both hosts.
Scoring, and keeping the whole live half, restores $16/16$.}
\label{tab:best}
\end{table}

Table~\ref{tab:best} holds the batch width fixed and compares two admission rules at $\tau_{\mathrm{dec}}$.
Ignoring the score and keeping one residual drops any-finished accuracy on both hosts.
Keeping every residual at or above $\tau_{\mathrm{dec}}$ restores it.
The savings in Table~\ref{tab:hosts} come from dropping loops, and at $\tau_{\mathrm{dec}}$ from dropping illegal text as well.
They do not come from collapsing the fan-out onto its winner.

\FloatBarrier
\section{Discussion}
\label{sec:discussion}

We presented a lightweight approach to efficient agentic inference.
By identifying and pruning abortable residuals during KV-cache encoding of the residual, and by continuing an uncertain residual only when a short prefix clears a second bar, our approach ensures that prefill and transfer are spent on live strategies rather than on text the search will abort.
Experiments show that \sys exhibits a host-invariant reduction in shared prefill, significantly lowering end-to-end time while preserving large-fan-out any-finished accuracy on GSM8K and MATH-500.
Collapsing either decision onto the other bar, or onto a single reserved slot, is what loses problems.
Our work highlights a practical split for enhancing the efficiency of LLMs in agentic serving, complementary to decode-time hosts, prefix caches, and prefill--decode disaggregation.

\section{Limitations}
\label{sec:limit}

The scores are surface tests: repetition, an illegal-math pattern, an answer mark, and character statistics.
They are not a learned reward model, and a residual that is wrong without being repetitive passes the prefill gate.
The evaluation mix is constructed so that half the fan-out is live strategy text and half is a loop or illegal string.
It is a controlled test of the two thresholds, not a trace of a production agent.
Any-finished accuracy can stay flat while winner accuracy moves, and at $k{=}4$ even $\tau_{\mathrm{pre}}$ drops one GSM8K problem.
The MATH-500 probe recovers full any-finished accuracy at $k{=}16$ and does not do so at $k{=}4$ or $k{=}8$.
The pruning study uses one model and temperature $0$, so the tables are single-run measurements.
The connector numbers in Appendix~\ref{app:xfer} are an accounting model, not a measured transfer time.
We do not explore training-based alternatives such as finetuning a pruner, since we target the setting without a heavy compute budget.
